\seccionTesis{Hipótesis}

\subseccionTesis{Hipótesis general}

Existe relación entre la condición del drenaje vial y el desempeño funcional del pavimento en el tramo EMP.LI105 Pedregal–Censo–Cumbray, Simbal, durante 2026.

\subseccionTesis{Hipótesis específicas}

\begin{listaTesis}
\item \textbf{H1.} La evacuación longitudinal presenta una relación directa con la condición superficial del pavimento.
\item \textbf{H2.} La evacuación longitudinal presenta una relación inversa con la rugosidad longitudinal del pavimento.
\item \textbf{H3.} La condición del drenaje transversal presenta una relación directa con la condición superficial del pavimento.
\item \textbf{H4.} La condición del drenaje transversal presenta una relación inversa con la rugosidad longitudinal del pavimento.
\end{listaTesis}

\seccionTesis{Variables}

\subseccionTesis{Variables de estudio}

La variable I es la condición del drenaje vial, expresada mediante continuidad longitudinal y abertura transversal libre, mientras que la variable II es el desempeño funcional del pavimento, caracterizado mediante el PCI y el IRI; ambas se registrarán en unidades estructura–sector, constituidas por una estructura transversal y su correspondiente sector de evaluación.

\subseccionTesis{Definición operacional}

La continuidad longitudinal libre se calculará mediante la expresión \(CL=100\,L_c/L_e\), donde \(L_c\) es la longitud del recorrido con conexión superficial continua hacia una descarga y \(L_e\) corresponde a la longitud total evaluada en ambos bordes del sector, expresadas en metros; cada interrupción se delimitará mediante progresivas y fotografía, contabilizando una sola vez los intervalos superpuestos.

La abertura transversal libre se determinará como \(CT=100\,\min(A_{le}/A_e,\,A_{ls}/A_s)\), donde \(A_e\) y \(A_s\) representan las áreas interiores de entrada y salida, respectivamente, y \(A_{le}\) y \(A_{ls}\) las correspondientes áreas libres, expresadas en metros cuadrados; de esta manera, el indicador representará el extremo con menor proporción de abertura disponible, sin equipararlo con la capacidad hidráulica de descarga.

Ambos indicadores del drenaje se expresarán en porcentajes comprendidos entre cero y cien, de modo que un valor mayor describirá una continuidad superior del recorrido o una mayor proporción de abertura libre. Por su parte, el PCI se calculará mediante el registro de deterioros, severidades, cantidades y valores deducidos conforme a ASTM D6433-18, mientras que el IRI corresponderá al promedio de los valores calculados para las dos huellas de circulación a partir de sus perfiles longitudinales, según ASTM E1926-08(2021) (ASTM International, 2018, 2021).

\subseccionTesis{Matriz de operacionalización de variables}

\begin{matrizOperacionalizacion}
\tituloMatrizOperacionalizacion[tab:matriz-operacionalizacion]{Operacionalización de la condición del drenaje y del desempeño funcional del pavimento}
\vspace{4pt}
\noindent\begin{tabularx}{\linewidth}{|
>{\raggedright\arraybackslash}p{2.45cm}|
>{\raggedright\arraybackslash}p{3.45cm}|
>{\raggedright\arraybackslash}p{4.2cm}|
>{\raggedright\arraybackslash}p{2.35cm}|
>{\raggedright\arraybackslash}X|
>{\raggedright\arraybackslash}p{1.8cm}|
>{\raggedright\arraybackslash}p{3.65cm}|}
\hline
\textbf{Variable} & \textbf{Definición conceptual} & \textbf{Definición operacional} & \textbf{Dimensión} & \textbf{Indicador} & \textbf{Tipo} & \textbf{Instrumento} \\ \hline
I. Condición del drenaje vial &
Condición física de recorridos y estructuras para conducir agua (Mays, 2019; Huang, 2004). &
Recorrido continuo respecto de la longitud evaluada: \(CL=100L_c/L_e\). &
Evacuación longitudinal &
Continuidad longitudinal libre (\%); longitud continua y evaluada (m). &
Cuantitativa &
Ficha geométrica, wincha, GPS y fotografías. \\ \hline
I. Condición del drenaje vial &
Condición física de recorridos y estructuras para conducir agua (Mays, 2019; Huang, 2004). &
Mínima proporción de área libre entre entrada y salida: \(CT=100\min(A_{le}/A_e,A_{ls}/A_s)\). &
Drenaje transversal &
Abertura libre (\%); áreas libres e interiores (m\textsuperscript{2}). &
Cuantitativa &
Ficha de estructuras, regla, wincha y fotografía escalada. \\ \hline
II. Desempeño funcional del pavimento &
Condición de la superficie respecto del servicio vial, observada por deterioros y regularidad (Shahin, 2005; Huang, 2004). &
Inspección de fallas, severidad, cantidad y deducciones según ASTM D6433-18. &
Condición superficial &
PCI (0–100). &
Cuantitativa &
Formato PCI y plantilla de cálculo. \\ \hline
II. Desempeño funcional del pavimento &
Condición de la superficie respecto del servicio vial, observada por deterioros y regularidad (Shahin, 2005; Huang, 2004). &
Promedio de IRI de las dos huellas, calculado desde el perfil longitudinal. &
Rugosidad longitudinal &
IRI (m/km). &
Cuantitativa &
Equipo de adquisición de perfil longitudinal y cálculo ASTM E1926-08(2021). \\ \hline
\end{tabularx}
\fuenteTabla{Elaboración propia, con base en la matriz de operacionalización del proyecto de tesis de Simbal, 2026.}
\end{matrizOperacionalizacion}

\seccionTesis{Objetivos}

\subseccionTesis{Objetivo general}

Determinar la relación entre la condición del drenaje vial y el desempeño funcional del pavimento en el tramo EMP.LI105 Pedregal–Censo–Cumbray, Simbal, durante 2026, mediante inspección georreferenciada, evaluación PCI, medición IRI y análisis correlacional por unidad espacial, para cuantificar su magnitud y dirección y orientar el seguimiento de la conservación vial.

\subseccionTesis{Objetivos específicos}

\begin{listaTesis}
\item Determinar la relación entre la evacuación longitudinal y la condición superficial, mediante inspección y análisis correlacional por sector, para orientar la conservación del pavimento.
\item Determinar la relación entre la evacuación longitudinal y la rugosidad longitudinal, mediante inspección, perfilado y análisis correlacional por sector, para orientar el seguimiento de la regularidad de rodadura.
\item Determinar la relación entre el drenaje transversal y la condición superficial, mediante inspección y análisis correlacional por sector, para orientar la conservación del pavimento junto a los pases de agua.
\item Determinar la relación entre el drenaje transversal y la rugosidad longitudinal, mediante inspección, perfilado y análisis correlacional por sector, para orientar el seguimiento de la regularidad de rodadura.
\end{listaTesis}
